\subsubsection{Ducros Marker}
\label{SEC:DucrosMarker}
\key*{}[marker!Ducros Marker]

%%%%%%%%%%%%%%%%%%%%%%%%%%%%%%%%%%%%%%%%%%%%%%%%%%%%%%%%%%%%%%%%%%%%%%%%%%%%%%%%

The Ducros marker uses a divergence based indicator together with the jump
marker to identify troubled cells.
%
First, the jump marker is computed using
\eref{EQN:JumpIndicator2,EQN:JumpMarker} and then the scaled jump marker,
$\scaledmarker_{\textrm{jump}}$, is evaluated using
\eref{EQN:ScaledResolutionMarker}.
%
Next, the divergence indicator is computed using the equation:
%
\begin{equation}
  \scaledmarker_{\textrm{div}} = \left(
                                 \frac{2 \left|\nabla\cdot\vec{V}\right|}{a}
                                 \cdot
                                 \frac{h}{\p[+1]{}}
                                 \right)^4
\end{equation}
%
where $h$ is the maximum length of a cell across all Cartesian coordinate
directions and $a$ is the speed of sound.
%
The value of $\scaledmarker_{\textrm{div}}$ for a cell is taken as the maximum
value from all solution points of the cell.
%
Finally, the scaled Ducros marker is computed as
%
\begin{equation}
  \scaledmarker_{\textrm{Ducros}} = \scaledmarker_{\textrm{jump}} \cdot
                                    \scaledmarker_{\textrm{div}}
\end{equation}

%%%%%%%%%%%%%%%%%%%%%%%%%%%%%%%%%%%%%%%%%%%%%%%%%%%%%%%%%%%%%%%%%%%%%%%%%%%%%%%%

Further information about the Ducros marker can be found in
\linecite{Koeplinger_2026}.

%%%%%%%%%%%%%%%%%%%%%%%%%%%%%%%%%%%%%%%%%%%%%%%%%%%%%%%%%%%%%%%%%%%%%%%%%%%%%%%%

\SeeAlso{
\inp[nmlLimiter::MarkerOption]{Marker\_Option},
\inp[nmlLimiter::MarkerVariable]{Marker\_Variable}
}


%%%%%%%%%%%%%%%%%%%%%%%%%%%%%%%%%%%%%%%%%%%%%%%%%%%%%%%%%%%%%%%%%%%%%%%%%%%%%%%%

